\section{Introduction}

A visitor enters a gallery and notices a small, unassuming device resting on a plinth beside a painting. Curious, she picks it up and raises it to her eyes. Through the viewfinder, the painting comes alive---brushstrokes lift from the canvas, the artist's preparatory sketches shimmer beneath the surface, and a constellation of historical references unfolds in the space around the frame. After a few seconds, she lowers the device and sets it back on the plinth. She moves on to the next room. But something pulls her back. She returns, picks up the viewer again, and this time notices details she missed---a hidden signature, a color shift that reveals a later restoration. She sets it down once more, turns to her companion, and says: ``You have to see this.''

This encounter---brief, voluntary, repeatable, and shareable---represents a form of mixed reality interaction that is both remarkably common and remarkably undertheorized. Museum viewfinders, exhibition stereoscopes, coin-operated binoculars at scenic overlooks, and increasingly, handheld AR experiences at cultural sites all follow the same pattern: a person picks up an optical device, holds it to their body for seconds, experiences an augmented view, and puts it down. The interaction is not a session to be endured but a \textit{glimpse} to be savored.

Yet the dominant approach in mixed reality design assumes sustained use. Head-mounted displays from research prototypes to commercial products---Meta Quest, Apple Vision Pro, Microsoft HoloLens---are engineered for wearing durations measured in minutes to hours. Their design priorities reflect this assumption: ergonomic weight distribution for extended comfort, adjustable straps for secure fit, ventilation to manage heat buildup, and software ecosystems that reward prolonged engagement. Donning and doffing these devices is treated as \textit{friction}---an overhead cost to be minimized so that the ``real'' interaction can begin \cite{knight2006comfort}. The implicit message is clear: the longer you wear it, the better.

This assumption creates a fundamental mismatch when MR encounters are transient. In museums, the average dwell time at an exhibit is 17--30 seconds \cite{vomlehn2001exhibiting}. In exhibitions and public demonstrations, visitors cycle through experiences in minutes. In educational settings, students pass AR devices around for brief, focused encounters. In all of these contexts, the interaction \textit{is} the wearing---not what happens during a sustained session, but the act of bringing a device to the body, glimpsing an augmented reality, and setting it down. Designing these encounters with devices optimized for sustained wear is like designing a doorway with the assumption that people will stand in it for an hour.

The mismatch extends beyond ergonomics into social dynamics. Research on the social acceptability of head-worn devices consistently demonstrates that wearables covering the face provoke discomfort, suspicion, and stigma among bystanders \cite{koelle2020social, koelle2017acceptability}. People wearing HMDs in public spaces appear disconnected, surveilling, or simply odd. These social costs are proportional to wearing duration: an hour in a VR headset at a caf\'{e} reads very differently from a two-second peek through a museum viewfinder. Yet social acceptability research has largely examined continuous wearing scenarios, leaving the temporal dimension unexplored.

We propose \textbf{temporal wearing} as an undertheorized design dimension that reframes this mismatch. Temporal wearing describes encounters where wearing a device for a brief, bounded duration---typically seconds---\textit{constitutes} the interaction rather than enabling it. In temporal wearing, donning and doffing are not overhead costs but productive acts: the don initiates a perceptual shift, the doff concludes it, and the cycle of donning and doffing structures the entire experience. The device is not a tool that happens to be worn; the wearing \textit{is} the tool.

This practice is not new. It has deep historical roots in optical culture. The lorgnette, popular from the 18th century onward, was a pair of lenses on a handle---raised to the eyes for a focused look, then lowered. Opera glasses served the same function in theatrical contexts. The monocle, the hand-held magnifying glass, the stereoscope, the tourist viewfinder: all are temporal wearing devices, designed not for sustained use but for repeated, brief encounters. What is new is not the practice but the \textit{formalization}. Despite centuries of temporal wearing and decades of wearable computing research, no framework has identified wearing duration as a primary design variable or theorized the don/doff cycle as an interaction in its own right. We argue that formalizing this long-standing practice has measurable design consequences: it reveals a design space that existing frameworks overlook and yields empirically grounded principles for a category of wearable interaction that has been practiced but never systematically studied.

Foregrounding wearing duration as a design dimension reveals relationships that existing frameworks miss. Where wearability research asks \textit{where} on the body a device should be placed \cite{gemperle1998wearability, zeagler2017where} and \textit{how} it should be attached, temporal wearing asks \textit{for how long} and \textit{how often}. Where transition research examines how users move between devices or interface states \cite{hubenschmid2025hybrid}, temporal wearing examines how users move between wearing and not-wearing as a \textit{cyclic} process---not a single transition from state A to state B, but a repeating rhythm of engagement and disengagement. And where calm technology envisions information moving between the center and periphery of attention \cite{weiser1996calm}, temporal wearing extends this to the device itself: the device moves between presence and absence, between body-coupled and body-decoupled states.

In this paper, we develop the temporal wearing concept and validate it through design and empirical study. Our contributions are:

\begin{itemize}
    \item[\textbf{C1}] \textbf{Temporal wearing as a design dimension} and a three-part \textbf{wearing duration taxonomy} (permanent, session, and temporal wear) that foregrounds an overlooked but consequential variable in wearable interaction design. Temporal wearing has been practiced for centuries but never formalized; we argue that formalizing it has measurable design consequences.
    \item[\textbf{C2}] A \textbf{design space} for temporal wearing devices defined by three measurable axes---body coupling, transition cost, and social legibility---populated with six prototype instantiations and historical precedents.
    \item[\textbf{C3}] \textbf{Two-phase empirical validation}: a controlled lab study (N=20) comparing temporal wearing prototypes against a session-wear baseline on transition time, re-engagement, social comfort, and cognitive load; and an in-situ deployment at a public exhibition revealing emergent temporal wearing behaviors.
    \item[\textbf{C4}] Six \textbf{design principles} for temporal wearing devices, grounded in empirical findings from both study phases, offering actionable guidance for designing MR encounters measured in seconds rather than sessions.
\end{itemize}
